\fancyhead[L]{\small\textsf{\textbf{MỤC 1.5}\quad Hàm số ngược và Logarit}}
\fancyhead[R]{\small\textsf{\textbf{59}}}

\noindent
\begin{minipage}[t]{0.26\textwidth}
    \vspace{0pt}
    \small
    \textbf{\textsf{Ký hiệu của Logarit}}\\
    \textsf{Hầu hết các sách giáo trình giải tích và các ngành khoa học, cũng như trên máy tính cầm tay, đều sử dụng ký hiệu $\ln x$ cho logarit tự nhiên và $\log x$ cho ``logarit thập phân'' (common logarithm), $\log_{10} x$. Tuy nhiên, trong các tài liệu toán học và khoa học nâng cao hơn cũng như trong các ngôn ngữ lập trình máy tính, ký hiệu $\log x$ thường được dùng để chỉ logarit tự nhiên.}
\end{minipage}%
\hfill
\begin{minipage}[t]{0.71\textwidth}
    \vspace{0pt}
    \noindent{\color{stewartcyan}\rule{2.5mm}{2.5mm}}\hspace{6pt}{\large\textbf{\textsf{Logarit tự nhiên}}}
    \vspace{6pt}

    Trong tất cả các cơ số $b$ khả dĩ cho logarit, chúng ta sẽ thấy trong Chương~3 rằng sự lựa chọn cơ số thuận tiện nhất chính là số $e$, đã được định nghĩa trong Mục~1.4. Logarit với cơ số $e$ được gọi là \textbf{logarit tự nhiên} (\textit{natural logarithm}) và có một ký hiệu đặc biệt:

    \vspace{6pt}
    \begin{center}
        \fcolorbox{stewartred}{white}{
            \parbox{0.5\textwidth}{
                \[
                    \log_e x = \ln x
                \]
            }
        }
    \end{center}
    \vspace{4pt}

    Nếu đặt $b = e$ và thay thế $\log_e$ bằng ``$\ln$'' trong (6) và (7), các tính chất định nghĩa hàm logarit tự nhiên trở thành

    \vspace{6pt}
    \noindent
    \begin{minipage}[c]{0.12\textwidth}
        \centering
        \fcolorbox{stewartred}{white}{\color{stewartred}\bfseries\small 8}
    \end{minipage}%
    \begin{minipage}[c]{0.88\textwidth}
        \fcolorbox{stewartred}{white}{
            \parbox{\dimexpr\textwidth-2\fboxsep-2\fboxrule}{
                \[
                    \ln x = y \iff e^y = x
                \]
            }
        }
    \end{minipage}

    \vspace{10pt}
    \noindent
    \begin{minipage}[c]{0.12\textwidth}
        \centering
        \fcolorbox{stewartred}{white}{\color{stewartred}\bfseries\small 9}
    \end{minipage}%
    \begin{minipage}[c]{0.88\textwidth}
        \fcolorbox{stewartred}{white}{
            \parbox{\dimexpr\textwidth-2\fboxsep-2\fboxrule}{
                \[\begin{aligned}
                    \ln(e^x) &= x \quad x \in \mathbb{R} \\[2pt]
                    e^{\ln x} &= x \quad x > 0
                \end{aligned}\]
            }
        }
    \end{minipage}

    \vspace{10pt}
    Đặc biệt, nếu chúng ta lấy $x = 1$, ta thu được

    \vspace{4pt}
    \begin{center}
        \fcolorbox{stewartred}{white}{
            \parbox{0.5\textwidth}{
                \[
                    \ln e = 1
                \]
            }
        }
    \end{center}
    \vspace{4pt}

    Kết hợp Tính chất~9 với Quy tắc~3 cho phép ta viết
    \[
        x^r = e^{\ln(x^r)} = e^{r \ln x} \qquad x > 0
    \]
    Như vậy, một lũy thừa của $x$ có thể được biểu diễn dưới dạng hàm mũ tương đương; chúng ta sẽ thấy điều này rất hữu ích trong các chương tiếp theo.

    \vspace{6pt}
    \noindent
    \begin{minipage}[c]{0.12\textwidth}
        \centering
        \fcolorbox{stewartred}{white}{\color{stewartred}\bfseries\small 10}
    \end{minipage}%
    \begin{minipage}[c]{0.88\textwidth}
        \fcolorbox{stewartred}{white}{
            \parbox{\dimexpr\textwidth-2\fboxsep-2\fboxrule}{
                \[
                    x^r = e^{r \ln x}
                \]
            }
        }
    \end{minipage}

    \vspace{14pt}
    \noindent\textbf{\textsf{\color{stewartcyan}VÍ DỤ 7}}\quad Tìm $x$ nếu $\ln x = 5$.

    \vspace{4pt}
    \noindent\textbf{\textsf{\color{stewartcyan}LỜI GIẢI 1}}\quad Từ (8), ta thấy rằng
    \[
        \ln x = 5 \quad \text{có nghĩa là} \quad e^5 = x
    \]
    Do đó $x = e^5$.

    (Nếu bạn gặp khó khăn khi làm quen với ký hiệu ``$\ln$'', hãy thay thế nó bằng $\log_e$. Khi đó phương trình trở thành $\log_e x = 5$; do đó, theo định nghĩa của logarit, $e^5 = x$.)

    \vspace{6pt}
    \noindent\textbf{\textsf{\color{stewartcyan}LỜI GIẢI 2}}\quad Bắt đầu với phương trình
    \[
        \ln x = 5
    \]
    và áp dụng hàm mũ cho cả hai vế của phương trình:
    \[
        e^{\ln x} = e^5
    \]
    Mặt khác, công thức triệt tiêu thứ hai trong (9) chỉ ra rằng $e^{\ln x} = x$. Do đó $x = e^5$.\hfill{\color{stewartcyan}$\blacksquare$}
\end{minipage}